Realistic reservoir simulation is known to be prohibitively expensive in terms of computation time when increasing the accuracy of the simulation or enlarging the model grid size. One method to address this issue is to parallelize the computation by dividing the model into several partitions and using multiple CPUs to compute the result using techniques such as MPI and multi-threading. Alternatively, GPUs are also a good candidate for accelerating computation due to their massively parallel architecture, which allows many floating point operations per second to be performed. Although Computational Flow Dynamics problems are complex and contain many computational parts that are challenging, the most difficult one is the execution of the numerical iterative solver that is used to solve the linear systems, which arise after discretizing the nonlinear system of equations that mathematically model the problem. The numerical iterative solver thus takes the most computational time and is challenging to solve efficiently due to the dependencies that exist in the model between cells. In this work, we evaluate the OPM Flow simulator and compare several state-of-the-art GPU solver libraries as well as custom-developed solutions for a BiCGStab solver using an ILU0 preconditioner and benchmark their performance against the default DUNE library implementation running on multiple CPU processors using MPI. The evaluated GPU software libraries include a manual linear solver in OpenCL and the integration of several third-party sparse linear algebra libraries, such as cuSparse, rocSparse, and amgcl. To perform our bench-marking, we use small, medium, and large use cases, starting with the public test case NORNE, which includes approximately 50k active cells, and ending with a large model that includes approximately 1 million active cells. We find that a GPU can accelerate a single dual-threaded MPI process up to 5.6 times and that it can compare with around 8 dual-threaded MPI processes.  

